\section{Описание}

В работе реализован словарь на основе дерева PATRICIA (Practical Algorithm To Retrieve Information Coded In Alphanumeric).
Такая структура представляет собой сжатое бинарное префиксное дерево, в котором в каждом узле хранится номер бита, по которому происходит ветвление.

Каждый узел содержит:
\begin{itemize}
    \item строковый ключ;
    \item связанное с ним значение типа \texttt{uint64\_t};
    \item индекс проверяемого бита;
    \item указатели на левого и правого потомка.
\end{itemize}

В качестве служебного элемента используется специальная вершина \texttt{head}, которая создаётся в конструкторе и не сохраняется во внешний файл. При сериализации ссылки на неё кодируются специальным значением \texttt{UINT32\_MAX}.

\subsection{Поиск}

Поиск выполняется от корня дерева вниз, пока индекс очередного узла строго возрастает. После завершения спуска проверяется совпадение найденного ключа с искомым словом.

\subsection{Вставка}

Перед вставкой слово переводится в нижний регистр. Если такого ключа в словаре ещё нет, вычисляется первый бит, на котором новый ключ отличается от уже найденного в дереве слова. В этот момент в дерево вставляется новая вершина, а её потомки распределяются в зависимости от значения соответствующего бита.

\subsection{Удаление}

Удаление оказалось самой сложной операцией. В корректной версии алгоритма пришлось отдельно обрабатывать случаи удаления терминальной вершины и удаления внутренней вершины с перепривязкой обратных ссылок. В ранней версии реализации здесь были ошибки, приводившие к висячим указателям.

\subsection{Сохранение и загрузка}

Для сохранения дерево обходится в фиксированном порядке и для каждого узла во внешний файл записываются:
\begin{enumerate}
    \item длина ключа;
    \item символы ключа;
    \item значение;
    \item индекс бита;
    \item индексы левого и правого потомков.
\end{enumerate}

При загрузке сначала создаётся временный словарь. После чтения всех узлов восстанавливаются связи между ними, затем выполняется проверка корректности структуры через поиск по сохранённым ключам. Только если вся операция прошла успешно, новый словарь подменяет текущий.

\section{Исходный код}

Основной алгоритм реализован в одном файле \texttt{main.cpp}. Центральной частью программы является класс \texttt{PATRICIA}, который инкапсулирует внутреннюю структуру дерева и все операции над словарём.

Внутри класса отдельно реализованы:
\begin{itemize}
    \item поиск слова по ключу;
    \item вставка нового слова;
    \item удаление существующего слова;
    \item сохранение дерева во внешний бинарный файл;
    \item загрузка дерева из файла.
\end{itemize}

Код программы можно разделить на две логические части. Первая часть отвечает за внутреннее устройство структуры данных: вычисление битов ключа, спуск по дереву, работу с вершинами и сериализацию. Вторая часть представляет собой консольный интерфейс, который считывает команды, вызывает соответствующие методы словаря и печатает результат в формате, заданном условием.

При сохранении и загрузке используются не только ключи и значения, но и параметры, необходимые для восстановления самой структуры дерева: индекс бита и номера потомков. За счёт этого дерево после загрузки восстанавливается в том же виде, в котором оно находилось до сохранения.

\section{Консоль}
\begin{alltt}
root$ ./main
+ a 1
OK
+ A 2
Exist
a
OK: 1
- a
OK
a
NoSuchWord
\end{alltt}
\pagebreak
